% ---- 摘要页（电子版首页） ----
\thispagestyle{abstractstyle}

\begin{center}
    \vspace*{0.5em}
    {\heiti\zihao{2} 基于时空联合估计与卡尔曼滤波的\\[0.3em]多源定位融合及任务规划}\\[0.8em]
    \rule{0.5\linewidth}{0.6pt}
\end{center}

\abstracttitle

移动机器人多源定位需同时处理采样异频、时间偏移和观测噪声。本文围绕两路异频（4\,Hz / 5\,Hz）定位数据，从时空对齐到任务规划逐层建模求解。

\textbf{问题一}（无噪声时间对齐）：两路数据仅存在开机时间差，对各路构建三次样条后以残差平方和为目标做粗-精两阶搜索。所得时间偏差 \textbf{$\Delta\tau = -198.43$\,s}，对齐残差降至 $10^{-11}$ 量级（机器精度），输出 10\,Hz 轨迹 7004 点。

\textbf{问题二}（含噪声联合估计）：噪声（$\sigma \approx 0.8$\,m）与固定空间偏差同时存在时，将 $(\Delta\tau, \Delta x, \Delta y)$ 三参数纳入同一目标函数联合求解，而非逐步估计。得 $\Delta\tau = -48.44$\,s，空间偏差 $(-3.44,\, 1.81)$\,m，残差 RMSE $= 1.56$\,m。常加速度卡尔曼滤波异步融合后输出 8513 点轨迹。

\textbf{问题三}（偏差存在性判定）：对真实数据同时计算 Wald 统计量（$W = 2.07$，$p = 0.36$）和 AIC 差值（$\Delta\text{AIC} = -2.81$），两项独立证据指向同一结论——\textbf{无显著系统偏差}，按零偏模型融合输出 3691 点。

\textbf{问题四}（任务调度优化）：在 SG 平滑轨迹上，对 36 个目标（18 射击 + 18 拍照）搜索满足距离/速度/加速度/准备时长约束的可行窗口，分别用贪心启发式和 CP-SAT 约束规划器求解。CP-SAT 在约 20\,s 内求得最优方案：\textbf{射击 10 次 + 拍照 15 次 = 25 个任务}。

\textbf{创新点：}（1）联合参数估计一步求解时间与空间偏差，避免分步估计的偏差传播；（2）Wald + AIC 双判据交叉验证，提高有限样本下的判别可信度；（3）贪心与 CP-SAT 互为校验，兼顾计算效率与全局最优性。

\keywords{\textbf{多源融合定位；联合参数估计；卡尔曼滤波；Wald 检验；贪心优化}}
